\newpage
% =================================================
\section{Exemplos de fluxos}
% =================================================

\subsection{Fluxos complementares e critérios de completude}
Os fluxos são considerados completos quando descrevem ator, pré-condição, entrada, validação, mudança de estado, efeitos colaterais e resultado. As operações críticas abaixo complementam os cenários de tela já descritos.

\subsubsection{Concessão do papel Bolsista}
O Chefe Geral seleciona um usuário com papel Aluno. O backend verifica a pré-condição ($Bolsista \implies Aluno$), rejeita a operação se o usuário não for Aluno ou gerar combinação proibida (ex.: Aluno Bolsista atuando como Gestor de Almoxarifado), cria o vínculo de papel na coleção de perfis, registra auditoria e sincroniza claims e só então anuncia a permissão disponível após renovação do token.

\subsubsection{Aceitação de convite}
O aluno apresenta um convite válido. O backend verifica e-mail normalizado, status, expiração e capacidade da turma. A aceitação é idempotente: o vínculo Aluno--Turma é criado uma única vez e o convite deixa de ser pendente. Usuário e papel Aluno são criados apenas quando ainda não existirem.

\subsubsection{Arquivamento e desarquivamento de turma}
O professor responsável altera o status da turma para Arquivada. Posts, comentários, histórico e vínculos com alunos não são apagados. O evento é auditado e notificações são emitidas quando aplicável. O desarquivamento restaura o status Ativo sem recriar a turma.

\subsubsection{Retirada de reagente}
O Gestor de Almoxarifado seleciona o frasco e o Professor ou Bolsista retirante. A transação relê o frasco, valida disponibilidade e escopo do gestor, registra o peso de saída, cria o empréstimo EM\_USO e muda a disponibilidade para EMPRESTADO. Se o frasco estiver sendo aberto pela primeira vez, sua validade é recalculada antes de aprovar. Se outro processo tiver retirado o frasco antes da confirmação, a transação falha sem criar empréstimo parcial.

\subsubsection{Devolução de reagente}
O Gestor informa o frasco e o peso de retorno. A transação verifica que o empréstimo está EM\_USO ou ATRASADO, recalcula consumo e atraso com dados do servidor (prazo encerrando às 23:59:59.999 do dia previsto), atualiza o empréstimo, grava o evento no histórico e devolve o frasco à disponibilidade adequada ou ao fluxo de descarte/quarentena se vencido ou esgotado.

\begin{explicacao}[Nota sobre o Login]
Quando não houver sessão válida, a etapa inicial é COM-01: o usuário acessa a tela de Login do LCQUI, insere suas credenciais (e-mail e senha) ou escolhe o "Login com Google". O sistema exibe um \textit{spinner} (ícone de carregamento giratório) no botão e carrega conta ativa e permissões antes de abrir a dashboard correspondente. Na mesma sessão, os fluxos seguintes não exigem repetir login.
\end{explicacao}

\subsection{Fluxos do usuário Gestor de Bem Patrimonial}

\subsubsection{Adicionando um bem patrimonial do LCQUI no sistema}
Partindo da tela de login, o Gestor acessa o sistema. Caso possua múltiplos papéis, ele clica no botão "Gestor de Bens Patrimoniais" no header. Na dashboard, clica no botão "Mais (+)" e seleciona "Novo Bem Patrimonial". Um modal se sobrepõe à tela. Ele preenche o número do patrimônio, pesquisa um \textit{Resumo\_Bem} existente, escolhe o Local (prédio, andar, sala) via menus \textit{dropdown} e anexa a URL da foto. Ao clicar em "Salvar", um \textit{spinner} roda no botão e, em seguida, um \textit{toast} (notificação flutuante) verde surge na tela: "Bem cadastrado com sucesso". O modal fecha e a grid central é recarregada mostrando o novo item.

\subsubsection{Baixando a lista de bens inservíveis para dar baixa no SEI}
Partindo da tela de login, o Gestor acessa a dashboard e clica no botão "Gerar Relatórios" na parte superior. O modal de relatórios abre. Ele seleciona o tipo "Bens Inservíveis", confirma e clica em "Gerar PDF". O botão muda para o estado \textit{loading}. Após alguns segundos, o navegador exibe a notificação de download concluído ou abre o arquivo em uma nova aba, pronto para ser anexado ao Sistema Eletrônico de Informações (SEI) da UENF.

\subsubsection{Baixando o relatório mensal Geral do Mês de Agosto}
Após o login, na tela principal, o Gestor clica em "Gerar Relatórios". No modal que se abre, seleciona a aba "Relatório Geral". Nos seletores de data, escolhe "Agosto" e "2026". Ao clicar em "Gerar", um círculo de progresso é exibido no centro do modal. Finalizado o processo backend, o modal exibe a mensagem "PDF pronto" e inicia o download automaticamente.

\subsubsection{Baixando relatório de bens materiais filtrado (Prédio P5, Andar 1)}
Partindo da tela de login, o Gestor abre o modal de relatórios e vai até a aba "Personalizado". Ele preenche as datas (15/06/2026 a 15/07/2026) e abre o filtro avançado. Marca a \textit{checkbox} do "Prédio P5" e digita "1" no andar. Clica em "Gerar PDF". Após a barra de carregamento preencher, o sistema entrega o documento formatado apenas com os equipamentos deste andar específico.

\subsubsection{Verificando requisições pendentes de professores}
Após efetuar o login, a dashboard do Gestor já apresenta no painel central um card de "Painel de Requisições". O botão de notificações (sino) está com uma bolinha vermelha. O Gestor clica na notificação e depois na requisição desejada. A tela faz uma transição para o modo \textit{Split-screen}. Na metade esquerda da tela, ele vê os dados estáticos (antigos) do patrimônio em texto cinza. Na metade direita, ele vê os dados novos sugeridos pelo professor (ex: alteração de sala) destacados em fundo amarelo suave, para chamar a atenção à mudança.

\subsubsection{Aprovando uma requisição pendente}
Dando continuidade ao fluxo de verificação (após login e abertura da tela de \textit{Split-screen}), o Gestor avalia que a alteração de sala pedida pelo professor está correta. Ele clica no botão verde massivo "Aprovar Alteração" no rodapé. Um ícone de \textit{check} ($\checkmark$) animado aparece no centro da tela. A requisição desaparece da fila de pendentes e um \textit{toast} de sucesso no canto superior direito confirma que o banco de dados principal foi atualizado.

\subsubsection{Alterando um bem patrimonial para: Dado baixa}
Partindo da tela de login, o Gestor usa a busca exata pela plaqueta canônica de um equipamento \textbf{Inservivel} que já passou por todo o processo burocrático no SEI. A tela de detalhes abre o rito próprio ``Registrar baixa'', não um editor genérico de status. O Gestor anexa o PDF no caminho autorizado; o backend valida objeto, dono/origem, tamanho, bytes \texttt{\%PDF-} e geração antes da transação. A confirmação relê autorização M9, status, versão e idempotência M7, grava comprovante imutável, evento histórico \texttt{baixa} e o status terminal \texttt{Ja\_dado\_baixa}. Falha de upload ou conflito preserva o estado anterior. Após sucesso o item sai do filtro ativo, continua consultável com seu histórico e a plaqueta não volta a ficar disponível.

\subsection{Fluxos do usuário Professor}

\subsubsection{Criando uma turma}
Após fazer login, o Professor está na sua dashboard. Ele clica no painel lateral esquerdo em "Turmas" e depois no botão flutuante "+ Nova Turma". Um modal simples é exibido solicitando Nome da Turma, Matéria, Ano, Semestre e Capacidade. Ele preenche "Química Orgânica I", seleciona a matéria correspondente, informa ano 2026, semestre 2 e capacidade 30, e clica em "Salvar". Um \textit{toast} de sucesso é disparado, um novo card desliza aparecendo no painel lateral, e o \textit{codigo\_turma} recém-criado pisca rapidamente na tela com um botão de "Copiar para a área de transferência".

\subsubsection{Adicionando Alunos em uma turma}
Com login feito, o Professor clica no card da turma recém-criada no painel lateral. A área central atualiza para a visão da disciplina. No header, ele clica no botão "Novo Aluno". No modal que surge, ele copia e cola os e-mails dos alunos separados por vírgula e clica em "Enviar Convites". Uma barra de carregamento é mostrada e o sistema apresenta o resultado por destinatário. O \textit{toast} só afirma envio de e-mail quando o envio externo foi acionado; a criação do registro de convite, por si só, é reportada como convite registrado, sem anunciar e-mail enviado. Falhas parciais permitem repetir somente os destinatários que faltaram.

\subsubsection{Criando um Post em uma turma}
Após login, o Professor acessa a disciplina clicando nela. No topo do Feed de publicações (área central), ele clica na caixa de texto "Comece um novo post...". O componente se expande. Ele preenche título e orientações da aula, clica no ícone de "Clipe de Papel" (anexo), o que abre o modal da sua biblioteca de roteiros em PDF. Ele seleciona o "Roteiro\_Prática\_3.pdf" e clica em "Publicar". A tela rola suavemente para baixo e o novo card da publicação surge com um efeito de \textit{fade-in}, exibindo o texto e a prévia do PDF anexado.

\subsubsection{Excluindo um Aluno de uma turma}
O Professor loga no sistema e navega até a turma desejada. Ele clica na aba "Membros" logo acima do feed. Uma lista com nome e foto dos alunos aparece. Ao lado do nome do aluno a ser removido, ele clica no ícone vermelho de lixeira. Um modal de dupla confirmação exige que ele clique em "Sim, remover aluno". A linha do aluno desaparece da lista com uma rápida animação de encolhimento (\textit{collapse}) e um \textit{toast} avisa "Aluno removido da turma".

\subsubsection{Respondendo um comentário de um aluno}
O Professor entra no sistema e nota uma bolinha vermelha piscando no sino de Notificações. Ele clica e vê "\textit{João comentou em um Roteiro}". Ao clicar na notificação, a tela rola automaticamente até a exata postagem no feed. Ele vê a dúvida do aluno, clica na caixa de texto abaixo dizendo "Escreva uma resposta..." e aperta Enter. O comentário dele surge instantaneamente logo abaixo do aluno, com seu nome identificado com a tag "[Professor]".

\subsubsection{Pesquisando Reagente Químico para checar disponibilidade}
Partindo da tela de login, o Professor clica no atalho "Reagentes" no painel esquerdo. O catálogo JSON de Resumos e Especificações é carregado no cliente (usando o cache local quando a versão publicada não mudou) e a busca textual ocorre localmente: na grande barra de pesquisa central, ele começa a digitar "Ácido Clori..." e a tabela filtra a cada caractere, sem filtrar por tecla no Firestore e sem exigir filtro prévio. Ele encontra o item, clica sobre a linha e um modal lateral (gaveta) desliza da direita mostrando em verde "\textit{Disponível: 3 frascos fechados, 1 aberto}".

\subsubsection{Pesquisando um bem patrimonial pelo Nome}
Após login, o Professor vai à área "Bens Patrimoniais" e digita "Microscópio" na busca. A tela é atualizada mostrando os equipamentos carregados que correspondem ao texto, após filtro obrigatório, com paginação (UI-04). Ele clica em um dos cards e vê a foto e a localização exata daquele Microscópio em tela cheia.

\subsubsection{Pesquisando um bem patrimonial pelo número de patrimônio}
Professor faz login, acessa a busca de "Bens Patrimoniais" e digita exatamente "987654". O sistema detecta a correspondência exata, pula a exibição da tabela genérica e direciona o professor direto para a página exclusiva de detalhes daquele patrimônio específico.

\subsubsection{Requisitando alteração de conservação (de Bom para Ruim)}
No passo anterior, ao estar na tela de detalhes do microscópio, o Professor, após logar, percebe que a lente está trincada. Ele clica no botão amarelo "Solicitar Edição". No formulário sobreposto, ele altera o campo \textit{dropdown} "Estado de Conservação" de "Bom" para "RUIM" (conservação); Inservivel é outro campo, de status, preenche o campo de texto justificando ("Lente principal rachada") e envia. O botão vira um \textit{spinner} rapidamente e um \textit{toast} avisa: "Requisição enviada ao Gestor de Patrimônio".

\subsubsection{Upload de Roteiro e compartilhamento com outros professores}
Professor acessa a dashboard após login, clica em "Gerenciar Roteiros" (no header central). O modal abre na aba "Meus". Ele clica no ícone de Nuvem para upload, arrasta um PDF do seu computador para a área pontilhada e aguarda a barra de progresso encher (100\%). Em seguida, no próprio arquivo recém-subido, ele clica em "Compartilhar". Uma lista de professores do sistema aparece, ele marca 3 caixinhas de seleção e clica no botão "Autorizar Acesso". Um alerta flutuante informa que o roteiro agora está visível nas bibliotecas dos colegas.

\subsection{Fluxos do usuário Chefe Geral}

\subsubsection{Criando um usuário Professor e usuário Gestor de Almoxarifado}
O Chefe Geral acessa a plataforma logando com suas credenciais. Na sua dashboard, acessa a aba "Professores" e clica no botão "+ Novo Professor". No modal, insere nome, e-mail institucional, laboratório e seleciona o Centro (ex: CCT). Ao salvar, um \textit{toast} avisa do envio do convite. Da mesma forma, indo para a aba "Almoxarifados", ele clica em "Mais (+)", depois "Novo Gestor de Almoxarifado", preenche o e-mail, e um novo card de gestor é populado silenciosamente na lista do sistema, notificando o sucesso da operação.

\subsection{Fluxos do usuário Aluno}

\subsubsection{Ingressando em uma Turma}
O Aluno acessa o site do LCQUI pela primeira vez, cria a senha (ou entra via Google) e faz o login. A dashboard dele é muito mais limpa e a tela central diz: "\textit{Você ainda não tem turmas}". Ele foca sua atenção no canto inferior esquerdo onde há o input "Código da Turma". Ele digita o \textit{código} de até 20 caracteres que o professor escreveu no quadro e clica em "Entrar". O botão gira um carregamento, o sistema verifica a capacidade máxima da turma no \textit{backend}, faz o link de tabelas, e instantaneamente um Card colorido com o nome da disciplina desliza para o centro da tela com o feedback: "\textit{Matrícula realizada com sucesso!}".

\subsubsection{Baixando um roteiro e realizando um comentário}
Logado, o aluno clica na disciplina recém-adicionada. O Feed da turma é carregado na tela. Ele rola a página e acha a postagem do professor com um arquivo. Ele clica no botão vermelho e muito visível de "Baixar PDF". O \textit{download} inicia via navegador. Logo abaixo do mesmo \textit{post}, na seção de interações, ele clica na caixa de texto vazia, digita "Professor, levaremos óculos de proteção?" e pressiona Enviar. O comentário surge na tela logo após uma animação suave de transição (\textit{fade-in}).

\subsection{Fluxos do Usuário Gestor de Almoxarifado}

\subsubsection{Adicionando um frasco fechado de um reagente existente}
Partindo do login, o Gestor de Almoxarifado acessa seu painel, clica no grupo "Estoque", digita o nome de um ácido na busca e clica em "Adicionar Frasco". Ele coloca o frasco lacrado em cima da balança USB (ou lê manualmente), digita no campo "Peso Total" (ex: 610g) e no "Volume Nominal" (ex: 500mL). Ao clicar em "Salvar", o sistema deriva a tara teórica (peso total menos o conteúdo nominal, sem caráter de medição física definitiva do recipiente vazio), registra no banco e o modal se fecha com uma notificação verde "Frasco Fechado registrado no almoxarifado".

\subsubsection{Adicionando um frasco já aberto de um novo reagente}
Após o login, o Gestor nota que o reagente ainda não existe no sistema. Na dashboard central, ele clica em "Novo Reagente". Preenche as propriedades químicas obrigatórias, com destaque em vermelho para o campo "Densidade" quando líquido. Ao confirmar, o modal avança automaticamente para o passo "Cadastrar Frascos Vinculados". Como o frasco físico já está aberto, ele seleciona a condição física inicial "Já Aberto", deixa o lote não informado (sem inventar lote) e informa apenas os dados reais disponíveis: o peso total atual medido na balança e, se o rótulo estiver legível, o conteúdo nominal. O sistema exibe "Saldo atual: desconhecido" e "Tara do recipiente: ainda não conhecida" e não pede tara conhecida, volume estimado nem massa estimada. Salva a operação e a tela é atualizada sem estimar saldo inicial; o frasco permanece com \texttt{saldo\_desconhecido = true} até que seja esgotado e pesado vazio.

\subsubsection{Realizando empréstimo (Retirada) de um frasco}
Com o login concluído, um Professor ou Bolsista chega no balcão. O Gestor de Almoxarifado clica rapidamente no enorme botão verde ``Registrar Retirada'' no centro da dashboard. Um modal abre solicitando a leitura do código do frasco ou pesquisa pelo nome. Localizado o item, o Gestor seleciona o usuário autorizado, informa o local de uso e pesa o frasco.

Antes da confirmação, a interface bloqueia frascos em quarentena ou pendentes de descarte. Se o frasco estiver vencido e seu uso vencido tiver sido autorizado, deve aparecer um alerta explícito: ``Este frasco venceu em DD/MM/AAAA. O uso vencido exige autorização e ciência conforme a finalidade; pesquisa requer TCR rastreável. Deseja confirmar o empréstimo mesmo assim?''. A confirmação é enviada ao backend, que repete a validação; o frontend não pode ignorar essa regra. Se a operação também registrar a primeira abertura do frasco, a validade efetiva é recalculada no servidor dentro da mesma transação.

\subsubsection{Devolvendo um frasco (Cálculo de Consumo)}
O professor devolve o vidro à bancada. O Gestor faz login (se necessário) e clica no grande botão ``Registrar Devolução''. Busca o item (que está na lista de pendentes). O sistema solicita ``Peso de Retorno'' e apresenta a confirmação explícita ``O frasco ficou vazio nesta devolução? [ ] Sim'', disponível tanto para frascos cadastrados como FECHADO quanto para os cadastrados como JA\_ABERTO. O sistema calcula o consumo e verifica a validade efetiva com o relógio do servidor (com prazo normalizado para as 23:59:59.999 do dia previsto). A confirmação do gestor é a autoridade do esgotamento, e não a comparação com a tara anterior. Ao confirmar que o frasco ficou vazio, o peso de retorno passa a ser a tara real: o frasco vai para VAZIO/INDISPONIVEL, o saldo vira conhecido como 0\,g (sólido) ou 0\,mL (líquido), registra-se FICOU\_VAZIO e o conteúdo nominal não é alterado. Se o gestor não confirmar o esgotamento, a diferença de peso é tratada como anomalia metrológica sujeita a justificativa e o frasco mantém seu estado. Não substituir o peso medido pela tara nem concluir uma devolução fictícia para contornar a validação. Se o frasco estiver vencido, ou tiver vencido enquanto permaneceu com o professor, a devolução permite registrar a decisão do gestor entre quarentena, pendência de descarte ou continuidade de uso excepcional conforme Q04; registrar descarte físico concluído é outra operação.

\subsubsection{Corrigindo uma retirada atribuída ao usuário errado}
Após o login, o Gestor identifica no histórico do frasco que a retirada foi vinculada ao Professor errado. Ele abre ``Corrigir operação'', que exibe a operação original, o efeito que será corrigido e um campo de motivo. O gestor lê o efeito e confirma. A correção é compensatória e auditável: o lançamento incorreto é preservado, a correção é vinculada à operação original com motivo e instante do servidor, e o vínculo do retirante é corrigido sem devolução fictícia, consumo fictício ou nova pesagem artificial. O peso físico de saída, o instante físico original e o local de uso permanecem registrados; o original não é apagado.

\subsubsection{Colocando um frasco em quarentena e decidindo sua saída}
Durante uma inspeção de bancada, o Gestor decide bloquear um frasco independentemente de cadastro, abertura, devolução ou reencontro. Ele usa ``COLOCAR EM QUARENTENA'', informa o motivo humano e confirma; o sistema relê o estado atual, valida o escopo do almoxarifado e grava \texttt{em\_quarentena = true} e \texttt{disponibilidade = INDISPONIVEL}, preservando vencimento, estado físico, saldo e validade. Na decisão posterior, o gestor escolhe ``VOLTAR A DISPONIVEL'' ou ``PENDENTE DE DESCARTE'', com motivo/laudo; a saída não revalida a validade nem altera o vencimento, e uma segunda decisão concorrente é rejeitada se a pré-condição já não existir.

\subsubsection{Extravio e reencontro de um frasco}
O extravio e o reencontro são fatos distintos de liberação. Os quatro fluxos abaixo obedecem às regras da subseção M5 da Seção~7.

\textbf{Fluxo A --- disponível $\rightarrow$ extravio $\rightarrow$ reencontro $\rightarrow$ quarentena $\rightarrow$ liberação.} O Gestor registra o extravio de um frasco disponível: ele assume \texttt{situacao\_localizacao = EXTRAVIADO} e \texttt{disponibilidade = INDISPONIVEL}, preservando o estado físico, \texttt{saldo}, \texttt{saldo\_desconhecido}, peso, tara, vencimento e \texttt{abertura\_historica\_desconhecida}; a autorização corrente de descarte é revogada, mantendo a decisão anterior no histórico; nenhum peso, consumo ou tara é inventado. Dias depois, o frasco é reencontrado; o gestor informa o estado físico constatado (\texttt{ABERTO}/\texttt{FECHADO}) e o sistema aplica \texttt{em\_quarentena = true} e \texttt{INDISPONIVEL}, sem marcar vencimento por causa do reencontro e sem reabrir empréstimo algum. O saldo permanece como estava. Após inspeção, o gestor decide ``VOLTAR A DISPONIVEL'' com motivo/laudo; como não há pendência metrológica, o estado físico é ABERTO/FECHADO e o frasco não está vencido sem autorização, a quarentena é encerrada e o frasco volta a DISPONIVEL sem revalidar validade.

\textbf{Fluxo B --- emprestado $\rightarrow$ extravio durante a custódia $\rightarrow$ encerramento extraordinário $\rightarrow$ reencontro $\rightarrow$ quarentena.} O frasco está EM\_USO com o Professor quando desaparece. O gestor registra o extravio com motivo humano: a localização vira EXTRAVIADO, o estado físico conhecido é preservado e a disponibilidade vira INDISPONIVEL e o empréstimo ativo é encerrado como \texttt{ENCERRADO\_EXTRAORDINARIO} com \texttt{EXTRAVIO\_SINISTRO}. Como não houve retorno físico, \texttt{peso\_retorno}, \texttt{peso\_retorno\_efetivo}, \texttt{medida\_utilizada} e \texttt{data\_devolucao\_efetuada} permanecem nulos, \texttt{consumo\_validado = false} e \texttt{peso\_saida} é preservado; \texttt{massa\_perda\_estimada\_g} só é gravada como estimativa se houver tara conhecida. Quando o frasco é encontrado, o gestor registra o reencontro e ele entra em quarentena; o empréstimo encerrado \textbf{não} é reaberto.

\textbf{Fluxo C --- reencontrado $\rightarrow$ quarentena $\rightarrow$ condição inadequada $\rightarrow$ pendência de descarte $\rightarrow$ descarte autorizado.} Após o reencontro, a inspeção conclui que o frasco não é reaproveitável. O gestor decide ``PENDENTE DE DESCARTE'' com motivo/laudo: a quarentena é encerrada, o frasco permanece INDISPONIVEL e é criada a autorização operacional estruturada de descarte (HQ-M2-008/B). Somente depois, com a autorização ativa, ``Descartar'' é permitido e o frasco vira DESCARTADO (terminal). Não existe caminho direto de quarentena para descarte.

\textbf{Fluxo D --- reencontrado com informação quantitativa insuficiente.} O frasco é reencontrado, mas não há tara conhecida nem medição confiável do conteúdo. A quarentena permanece; o saldo continua desconhecido e \textbf{não} é fabricado. Uma pesagem de rotina posterior pode ser registrada para observar o estado, mas não resolve automaticamente a quarentena nem a pendência quantitativa. Enquanto a informação for insuficiente, o frasco não volta a DISPONIVEL e a decisão de destino aguarda subsídio; qualquer determinação quantitativa definitiva pertence à fronteira M6 (Q06/tara), sem inventar consumo, tara ou saldo. Se o gestor optar por não reaproveitar, aplica-se o Fluxo~C.

\textbf{Casos de regressão documental.} O reencontro valida o estado constatado contra o último estado conhecido antes do extravio:
\begin{itemize}
    \item \textbf{Caso 1 (válido).} \texttt{FECHADO} $\rightarrow$ \texttt{EXTRAVIADO} $\rightarrow$ \texttt{FECHADO} $\rightarrow$ quarentena, somente se nunca houve abertura conhecida.
    \item \textbf{Caso 2 (válido).} \texttt{FECHADO} $\rightarrow$ \texttt{EXTRAVIADO} $\rightarrow$ \texttt{ABERTO} $\rightarrow$ quarentena, marcando \texttt{abertura\_historica\_desconhecida = true}, \texttt{data\_abertura = null}, \texttt{validade\_desconhecida = true} e \texttt{validade\_efetiva = null}.
    \item \textbf{Caso 3 (rejeitado).} \texttt{ABERTO} $\rightarrow$ \texttt{EXTRAVIADO} $\rightarrow$ \texttt{FECHADO}: um frasco já aberto não pode ser ``ressuscitado'' como fechado.
    \item \textbf{Caso 4 (válido).} \texttt{VAZIO} $\rightarrow$ \texttt{EXTRAVIADO} $\rightarrow$ \texttt{VAZIO} $\rightarrow$ quarentena $\rightarrow$ \texttt{PENDENTE\_DE\_DESCARTE}: restauração sem nova medição, seguida de pendência de descarte.
    \item \textbf{Caso 5 (rejeitado).} \texttt{VAZIO} $\rightarrow$ \texttt{EXTRAVIADO} $\rightarrow$ \texttt{FECHADO}.
    \item \textbf{Caso 6 (válido).} \texttt{QUEBRADO} $\rightarrow$ \texttt{EXTRAVIADO} $\rightarrow$ \texttt{QUEBRADO} $\rightarrow$ quarentena $\rightarrow$ \texttt{PENDENTE\_DE\_DESCARTE}: quebra física constatada, sem recuperação operacional.
    \item \textbf{Caso 7 (rejeitado).} \texttt{QUEBRADO} $\rightarrow$ \texttt{EXTRAVIADO} $\rightarrow$ \texttt{ABERTO}.
    \item \textbf{Caso 8 (sem fabricação).} Frasco anteriormente \texttt{ABERTO} reencontrado aparentemente sem conteúdo: M5 \textbf{não} inventa \texttt{VAZIO}, \texttt{saldo = 0} nem nova tara; mantém a quarentena, registra a observação e remete a confirmação quantitativa a M6.
\end{itemize}

\textbf{Casos de regressão documental de M8 (estoque, escassez, cache e notificações).} A aptidão usa o predicado único \texttt{frascoAptoParaUso} da Seção~\ref{sec:m8-estoque-escassez-notificacoes}.

\emph{Estoque atual.}
\begin{itemize}
    \item \texttt{FECHADO + LOCALIZADO + DISPONIVEL}: apto.
    \item \texttt{ABERTO + LOCALIZADO + DISPONIVEL}: apto.
    \item \texttt{ABERTO + EXTRAVIADO}: não apto, apesar do estado físico aberto.
    \item \texttt{QUEBRADO}, \texttt{VAZIO}, \texttt{DESCARTADO}, \texttt{EMPRESTADO}: não aptos.
    \item \texttt{em\_quarentena = true}: não apto, mesmo \texttt{ABERTO}/\texttt{FECHADO}.
    \item Frasco vencido com \texttt{uso\_vencido\_autorizado = true}: apto (validade compatível); sem autorização: não apto.
    \item Saldo desconhecido: a contagem de frascos é possível conforme a métrica, mas a soma quantitativa é segregada e \textbf{nunca} tratada como zero conhecido.
    \item Massa (g) e volume (mL) nunca são somados no mesmo total.
\end{itemize}

\emph{Escassez.}
\begin{itemize}
    \item \texttt{qtdAptos < limiar}: escassez (emite alerta quando configurado).
    \item \texttt{qtdAptos = limiar}: \textbf{não} há escassez.
    \item \texttt{qtdAptos > limiar}: não há escassez.
    \item \texttt{ativo = false}: não avaliar.
    \item \texttt{notificacao\_ativa = false}: escassez pode existir conceitualmente, sem alerta automático.
    \item \texttt{Almoxarifado.ativo = false}: não executar avaliação automática.
    \item Frascos extraviados, quebrados, vazios, descartados, emprestados ou em quarentena não inflam \texttt{qtdAptos}.
\end{itemize}

\emph{Cache.}
\begin{itemize}
    \item \texttt{cache hit} válido e autorizado: responde dentro da validade restante.
    \item \texttt{cache hit} sem autorização: rejeita.
    \item \texttt{cache hit} acima do rate limit: rejeita ($5$/min/UID).
    \item \texttt{cache miss} ou inválido: recalcula sob demanda.
    \item Geração muda durante o cálculo: não publica resultado obsoleto.
    \item Duas mudanças consecutivas: responde indisponibilidade temporária, não dado obsoleto.
    \item TTL físico atrasado: não torna o cache semanticamente válido.
    \item $30$ segundos sem acesso: nenhum recálculo automático (não há scheduler).
    \item Mutação de domínio atualiza geração e marca \texttt{valido = false} na \textbf{mesma} transação, sem recalcular a agregação.
    \item Correção exclusivamente histórica não invalida estoque atual.
\end{itemize}

\emph{Notificações.}
\begin{itemize}
    \item Primeira escassez do dia: cria \texttt{ESCASSEZ\_ESTOQUE} com docId determinístico.
    \item \textit{Retry} do job no mesmo dia: não duplica (\texttt{create}/\texttt{ALREADY\_EXISTS}).
    \item Mesma configuração no dia civil seguinte: nova notificação permitida.
    \item Erro diferente de \texttt{ALREADY\_EXISTS}: \textbf{não} engolido, propaga.
    \item Gestor sem vínculo ao almoxarifado: não recebe, mesmo com o papel global.
    \item ``Limpar tudo'': marca como lida, preserva histórico, sem \texttt{DELETE}.
\end{itemize}

\emph{Interação M5/M6.}
\begin{itemize}
    \item Frasco extraviado: fora de \texttt{qtdAptos}.
    \item Reencontrado em quarentena: fora de \texttt{qtdAptos}.
    \item Quebrado: fora de \texttt{qtdAptos}.
    \item Pendência metrológica bloqueante (\texttt{DEVOLVIDO\_COM\_ANOMALIA} com \texttt{consumo\_validado = false}): fora de \texttt{qtdAptos}.
    \item Resolução metrológica que mantém quarentena: continua fora.
    \item Liberação M5 válida (\texttt{VOLTAR\_A\_DISPONIVEL}): volta a depender das demais condições de aptidão.
    \item Consumo pendente não é contado como consumo zero nem como frasco apto adicional.
\end{itemize}

O sistema renderiza na tela em negrito e com fundo destacado: \textbf{``Volume utilizado neste empréstimo: 4.5 mL''} ou \textbf{``Peso utilizado nesse empréstimo: 22 g''}. Ao clicar em ``Finalizar'', um sinal sonoro/visual confirma que a transação de estoque, consumo, histórico e auditoria foram concluídas.


\subsection{Fluxos M11 --- Turmas, matrícula e convites}
Estes exemplos são o contrato normativo da Seção~\ref{sec:regras-turmas-m11}.
Eles são ilustrativos e não atestam a implementação.

\subsubsection{Corrida pela última vaga}
Turma com \texttt{qtd\_alunos = capacidade - 1}. Dois alunos aceitam ingressar por
código no mesmo instante. Cada transação lê \texttt{qtd\_alunos} e valida
\texttt{qtd\_alunos < capacidade}; o Firestore reexecuta as transações em
conflito. A primeira commita e grava vínculo, espelho, evento e
\texttt{qtd\_alunos + 1}. A segunda, já com \texttt{qtd\_alunos = capacidade},
recebe \texttt{failed-precondition} e não cria vínculo, espelho, evento nem
contador parcial.

\subsubsection{Dois convites simultâneos}
Para o mesmo par (e-mail normalizado, turma), ou para o mesmo e-mail com
\texttt{id\_turma} \texttt{NULL}, duas criações concorrentes disputam a chave
determinística em \texttt{Chaves\_Unicas}. A primeira cria o convite
\texttt{pendente}; a segunda, na reexecução, encontra a chave ocupada e é
rejeitada sem criar convite duplicado. Turmas diferentes para o mesmo e-mail
são permitidas. Reenviar para um convite pendente substitui hash e expiração no
mesmo documento em vez de criar outro.

\subsubsection{Replay de aceite}
O aluno aceita um convite e o vínculo é criado. Ao repetir a mesma operação
M7, o backend encontra o receipt e devolve o resultado persistido: o convite
permanece aceito, o vínculo já existe, o espelho está coerente e
\texttt{qtd\_alunos} não é incrementado de novo. Convite consumido por outro
e-mail ou já expirado é negado com erro específico, sem efeito parcial.

\subsubsection{Remoção concorrente}
Professor dono e Chefe (Q13) removem o mesmo aluno simultaneamente. A primeira
transação apaga o vínculo canônico, o espelho, grava \texttt{exclusao\_aluno}
com \texttt{removido\_por} e decrementa \texttt{qtd\_alunos}. A segunda
encontra o vínculo ausente, trata o caso como idempotente e não decrementa o
contador de novo, jamais abaixo de zero. Dois retries de remoção pela mesma
origem seguem a mesma regra. Posts e comentários anteriores permanecem e nova
entrada por código é bloqueada, exigindo convite explícito.

\subsubsection{Arquivamento e desarquivamento}
O professor dono arquiva uma turma com alunos. A transação grava
\texttt{status = Arquivada} e incrementa \texttt{versao}, sem apagar membros,
posts, comentários, histórico nem alterar \texttt{codigo\_turma}. Enquanto
arquivada, o backend e as Rules rejeitam novo ingresso, novo convite, novo post
e novo comentário (Q08); os espelhos dos alunos recebem \texttt{status} por
fan-out idempotente, e espelho atrasado nunca concede escrita. O
desarquivamento restaura \texttt{Ativo} com o mesmo ID e os mesmos vínculos.

\subsubsection{Convite global}
O professor convida um e-mail sem turma, mediante confirmação explícita. O
convite tem \texttt{id\_turma} \texttt{NULL} e chave global própria. Ao aceitar,
o backend garante \texttt{Usuarios} e o papel Aluno e consome o convite, mas
\textbf{não} cria matrícula nem incrementa \texttt{qtd\_alunos}; o aluno entra
em alguma turma depois, por código ou por outro convite. Matrícula é
obrigatória e única na identidade final \texttt{Aluno}, opcional e não única no
convite.

\subsection{Fluxos M12.1 --- posts, comentários, edição e moderação}
Estes exemplos são o contrato normativo da
Seção~\ref{sec:regras-posts-comentarios-m12-1}. São ilustrativos e não atestam a
implementação.

\subsubsection{Publicar Post com e sem roteiro}
O professor dono abre a turma Ativo e publica um Post. Sem roteiro, a transação
relê o vínculo de ownership e o estado da turma, grava titulo/descricao e o
\texttt{idOperacao} M7 e emite a notificação \texttt{POST} como efeito
delimitado. Com roteiro, o backend valida o acesso atual ao roteiro, grava o
snapshot imutável \texttt{roteiro\_anexo} e só então publica; se o acesso não
puder ser validado, o Post não é criado (fronteira M12.2).

\subsubsection{Editar Post e remover da apresentação}
O professor dono edita o Post: a transação relê ownership, grava o evento em
\texttt{Historico\_Posts\_Turma} com valores antigo/novo, operador e timestamp,
marca \texttt{editado}/\texttt{editado\_em} e é idempotente M7; uma segunda
edição concorrente reexecuta a transação e preserva a versão anterior no
histórico, sem sobrescrita silenciosa. Remover da apresentação exige motivo e
não apaga o documento nem o histórico: o item some do feed de colegas e
permanece acessível ao professor dono e à auditoria.

\subsubsection{Moderar Post por Q13}
O Chefe identifica conteúdo impróprio em turma alheia e remove o Post da
apresentação com motivo. A transação exige turma Ativo, grava
\texttt{removido\_da\_apresentacao}, \texttt{motivo\_remocao},
\texttt{removido\_por}/\texttt{removido\_em}, registra o evento em
\texttt{Historico\_Posts\_Turma} (\texttt{tipo = moderacao}) e a auditoria,
sem assumir a autoria. Em turma arquivada a intervenção é negada.

\subsubsection{Comentar, editar e moderar comentário}
Aluno/Bolsista com vínculo canônico atual comenta; o texto é validado (1--2\,000,
sem HTML), escapado na renderização e gravado com \texttt{idOperacao} M7. O
autor edita o próprio comentário: grava \texttt{Historico\_Comentario},
marca \texttt{editado}/\texttt{editado\_em} e não desfaz moderação. O professor
dono (ou o Chefe por Q13) modera com motivo: \texttt{moderado = true},
\texttt{motivo\_moderacao}, \texttt{moderado\_por}/\texttt{moderado\_em}
atômicos, sem apagar o \texttt{texto}. Na listagem autorizada, o autor vê o
original marcado, o colega vê o aviso institucional e o professor dono/Chefe
veem original e histórico.

\subsubsection{Turma arquivada, remoção e replay}
Com a turma \texttt{Arquivada}, criar/editar/remover/moderar Post e
comentar/editar/moderar são negados, inclusive para o Chefe. Após a remoção de
um aluno, a leitura da turma e dos comentários pelo removido é negada, mas os
comentários anteriores permanecem para os demais. O replay da mesma operação M7
devolve o receipt e não duplica comentário, histórico ou notificação; reuso
incompatível de identidade de comando é rejeitado.

\subsubsection{Notificação sem acesso}
Uma notificação \texttt{POST}/\texttt{COMENTARIO} não carrega conteúdo
protegido. Ao clicar, o cliente revalida o acesso corrente; se o aluno foi
removido, a turma foi arquivada ou o item foi moderado, a navegação falha sem
exibir o original a terceiros; a notificação não autoriza leitura por si.

\subsection{Fluxos M12.2 --- roteiros, compartilhamento e download}
Estes exemplos são o contrato normativo da
Seção~\ref{sec:regras-roteiros-m12-2}. São ilustrativos e não atestam a
implementação.

\subsubsection{Upload e publicação de roteiro}
O professor ativo envia um PDF. O backend valida App Check/Auth/papel, tamanho
estritamente menor que 15 MiB, bytes \texttt{\%PDF-}, caminho e titularidade, e
fixa a geração validada antes de tornar o Roteiro publicável. Cancelar ou falhar
não cria registro utilizável; objeto órfão é reconciliado sem apagar objeto
referenciado. Um Roteiro não publicável não aparece em ``Meus''.

\subsubsection{Compartilhar e revogar (Q09)}
O proprietário seleciona professores ativos e autoriza o acesso; a relação
única por (roteiro, professor) é criada de forma idempotente e emite a
notificação mínima \texttt{ROTEIRO\_COMPARTILHADO}. Revogar remove o UID do
array no servidor: o destinatário deixa de ver o roteiro em ``Compartilhados
comigo'' e não pode criar novas associações, mas os Posts já publicados e seus
snapshots permanecem (Q09).

\subsubsection{Download por professor}
O proprietário baixa sempre. O destinatário baixa enquanto o compartilhamento
vigorar. Um professor terceiro não baixa. O servidor emite URL curta com a
identidade e a geração do objeto; a URL já emitida não é revogável antes de
expirar.

\subsubsection{Download por aluno, ex-aluno e turma arquivada}
O aluno com vínculo canônico atual baixa o anexo de um Post acessível, mesmo com
a turma arquivada (somente leitura). Um ex-aluno, ainda que a claim esteja
atualizada, não baixa nem lê o anexo; nenhum cache, erro ou notificação vaza o
conteúdo. Post removido da apresentação e roteiro revogado fora de Post
acessível não liberam download.

\subsubsection{Anexar, editar e desvincular Post}
O professor dono publica um Post com roteiro; o backend exige
\texttt{acessoRoteiroValidado} no commit e grava o snapshot imutável
\texttt{roteiro\_anexo} com a geração. Editar o Post mantendo ou trocando o
anexo exige acesso atual; sem ele, a edição falha fechado e o professor deve
remover o anexo. Desvincular não exige acesso e registra a versão anterior no
histórico de M12.1. O snapshot rotula a publicação, mas a leitura revalida o
acesso atual e nunca usa o snapshot como autorização permanente.

\subsection{Cobertura operacional por papel e ação}
\label{sec:fluxos-completos}
Cada fluxo parte de sessão ativa validada por COM-01 e usa os campos/estados do contrato UI referenciado na seção 8. Os nomes de coleções e os efeitos descritos seguem o dicionário da seção 5.9. Não repetir login entre ações da mesma sessão. A coluna de resultado inclui persistência e falhas relevantes.

\begin{regra}[M9 --- pré-condição comum dos fluxos]
Sessão autenticada não é autorização. Antes de cada leitura sensível ou
mutação, resolver a decisão atual da Seção~7/M9; para efeito crítico, ler
Usuário ativo, papel, escopo/ownership, recurso e estado de domínio na mesma
transação. Fluxo iniciado antes de revogação não pode gravar depois dela. Retry
M7 concluído apenas retorna o resultado já persistido ao mesmo UID.
\end{regra}

\subsubsection{COM-01 --- Todos: Login, recuperar senha, perfil, sair e alternar papel}
\textbf{Contrato:} UI-01. Login ou foto do header; preencher credenciais/perfil ou escolher papel autorizado. Auth e Usuarios determinam sessão; renovar permissões antes da dashboard. Cancelamento preserva sessão anterior; revogação limpa dados e redireciona.

\subsubsection{CHE-01 --- Chefe: Criar usuário e atribuir Bolsista ou outro papel}
\textbf{Contrato:} UI-02. Aba Pessoas > Novo > selecionar existente/convidar > preencher dados do papel > revisar > confirmar. Validar combinações e identidade, gravar perfil e auditoria, sincronizar Auth; lista atualiza após confirmação, falha de claims aparece como pendência.

\subsubsection{CHE-02 --- Chefe: Revogar papel ou autorrevogar chefia}
\textbf{Contrato:} UI-02. Pessoa > Papéis > Revogar > justificar > revisar conjunto restante. Transação verifica último chefe/gestor e vínculos; rejeição mantém papéis e audita tentativa. Último papel removido desativa conta e encerra acesso.

\subsubsection{CHE-03 --- Chefe: Criar/ativar almoxarifado e vincular gestores}
\textbf{Contrato:} UI-03. Almoxarifados > Mais > Novo > nome/descrição/local/gestores > Salvar. Criar local antes se necessário. Persistir almoxarifado e vínculos; ativo exige gestor. Gerenciar gestores permite trocar responsáveis sem deixar almoxarifado ativo órfão.

\subsubsection{CHE-04 --- Chefe: Consultar pessoas, turmas e atividade mensal}
\textbf{Contrato:} UI-02. Selecionar aba, filtro e pessoa > Atividade/Turmas > período. Ler perfis mínimos e agregados por autor/mês; abrir eventos paginados. Sem atividade exibe zero registros, sem confundir com erro de autorização.

\subsubsection{CHE-05 --- Chefe: Operar catálogo, patrimônio, matéria, etiquetas e relatórios}
\textbf{Contrato:} UI-03 a UI-09 e UI-12. Usar as mesmas entradas dos fluxos ALM/PAT/MAT abaixo com escopo global. Poder administrativo não dispensa campos, comprovante de baixa, destino de vencidos ou auditoria. Gestão de turmas/roteiros de terceiros requer Q13; não simular propriedade.

\subsubsection{MAT-01 --- Professor/Chefe: Cadastrar e editar matéria}
\textbf{Contrato:} UI-03. Mais > Nova Matéria ou detalhe > Editar > nome/código > confirmar. Reservar código único e persistir Materia; atualizar projeções e manter vínculos por ID. Código em uso retorna erro no campo.

\subsubsection{PRO-01 --- Professor: Criar turma e copiar código}
\textbf{Contrato:} UI-10. Turmas > Nova > nome/matéria/ano/semestre/capacidade > Salvar. Criar Turma e código único; selecionar turma e Copiar. Falha da área de transferência não recria turma.

\subsubsection{PRO-02 --- Professor: Arquivar e desarquivar turma}
\textbf{Contrato:} UI-10. Turma própria > Arquivar > confirmar; ou Turmas Arquivadas > selecionar > Desarquivar. Alterar status mantendo posts/membros; propagar espelhos e notificações. Q08: turmas arquivadas ficam somente em leitura.

\subsubsection{PRO-03 --- Professor: Convidar aluno com ou sem turma}
\textbf{Contrato:} UI-10. Mais/Novo Aluno > e-mails e matrícula opcional > turma opcional > confirmar ausência de turma ou justificar exceção de capacidade > Enviar. Persistir convites sem conceder papel diferente de Aluno; resultado por destinatário. Aceitação é ALU-01.

\subsubsection{PRO-04 --- Professor: Gerenciar membros e remover aluno}
\textbf{Contrato:} UI-10. Turma própria > Membros > selecionar > Remover > confirmar. O servidor remove o vínculo principal (\texttt{Turma/\{id\}/Alunos}) e também propaga a exclusão no espelho inverso (\texttt{Usuarios/\{uid\}/Turmas}) na mesma transação (N-20), atualiza contador e grava histórico/auditoria; comentários permanecem. Nova entrada por código bloqueada, retorno exige convite.

\subsubsection{PRO-05 --- Professor: Cadastrar, visualizar e baixar roteiro}
\textbf{Contrato:} UI-11. Mais > Novo Roteiro ou Gerenciar Roteiros > upload > nome/descrição/PDF > registrar. Servidor valida objeto e autoria. Biblioteca > Visualizar/Baixar revalida acesso e entrega arquivo; upload falho não cria roteiro utilizável.

\subsubsection{PRO-06 --- Professor: Compartilhar, acompanhar e revogar compartilhamento}
\textbf{Contrato:} UI-11. Biblioteca > Meus > Compartilhar > professores > confirmar. Criar vínculo único e notificação; Eu compartilhei mostra destinatários. Revogar exige confirmação e tratamento de posts existentes conforme Q09; recebido não autoriza redistribuição pelo destinatário.

\subsubsection{PRO-07 --- Professor: Publicar/editar post e vincular roteiro próprio ou recebido}
\textbf{Contrato:} UI-11. Turma própria > compositor/Editar > título/descrição > clipe > escolher roteiro ou cadastrar novo > Publicar. Revalidar acesso ao PDF e turma; persistir Post/histórico, notificar membros e atualizar feed. Falha mantém rascunho na sessão.

\subsubsection{PRO-08 --- Professor: Responder comentário e navegar por notificação}
\textbf{Contrato:} UI-11 e UI-12. Sino > aviso > post alvo > escrever comentário > Enviar. Revalidar propriedade da turma, persistir Comentario e notificar interessados; mostrar resposta no mesmo post, sem pressupor threads hierárquicas.

\subsubsection{PRO-09 --- Professor: Consultar reagentes, receber empréstimo e acompanhar consumo}
\textbf{Contrato:} UI-04, UI-07 e UI-12. Visualizar > Reagentes > filtros > detalhe > procurar gestor para retirada ALM-04. Meus Reagentes mostra empréstimos próprios e consumo histórico; somente gestor confirma movimentação. Consulta não cria reserva de estoque.

\subsubsection{PRO-10 --- Professor: Pesquisar patrimônio e requisitar adição/edição}
\textbf{Contrato:} UI-04 e UI-09. Bens > filtros/número > detalhe > Solicitar Edição, ou Mais > Pedido de Adição. Preencher proposta/motivo, criar Local se ausente, enviar. Lock e requisição são atômicos; Minhas Requisições acompanha resposta PAT-03.

\subsubsection{ALU-01 --- Aluno: Aceitar convite ou ingressar por código}
\textbf{Contrato:} UI-10. Convite > autenticar e verificar e-mail > completar nome/matrícula > Aceitar; ou painel > código > Entrar. Servidor verifica validade, combinações, turma e capacidade/exceção; cria vínculo/espelho/histórico uma vez. Convite inválido ou lotação não cria matrícula parcial.

\subsubsection{ALU-02 --- Aluno: Visualizar turmas, colegas, posts e baixar roteiro}
\textbf{Contrato:} UI-10 e UI-11. Minhas Turmas > turma > Feed/Membros > post > PDF. Ler somente turma com vínculo atual, colegas com dados mínimos e roteiro vinculado. Remoção invalida consultas e acesso ao arquivo; não confiar em espelho/cache antigo.

\subsubsection{ALU-03 --- Aluno: Comentar e ler notificações}
\textbf{Contrato:} UI-11 e UI-12. Post de turma matriculada > texto > Enviar; ou Sino > aviso > post. Validar vínculo e texto, persistir comentário e apresentar confirmação; Limpar tudo marca avisos lidos sem apagar fatos.

\subsubsection{ALU-04 --- Aluno: Consultar estoque sem retirar}
\textbf{Contrato:} UI-04. Visualizar > Reagentes > busca local (catálogo JSON) > detalhes de catálogo. Leitura conforme matriz da seção 3, sem identidade de retirantes nem botões de movimentação; concessão de Bolsista não pode ser feita pelo próprio aluno.

\subsubsection{BOL-01 --- Bolsista: Ações acadêmicas, consulta e retirada/devolução como destinatário}
\textbf{Contrato:} UI-04, UI-07, UI-10 a UI-12. Executar ALU-01 a ALU-04 pelo papel Aluno. Consultar reagente, comparecer ao balcão e informar finalidade/local/prazo ao gestor; ALM-04/05 registram saída/retorno. Acompanhar empréstimos próprios; nunca executar mutação de gestor nem acumular Gestor de Almoxarifado.

\subsubsection{ALM-01 --- Gestor Almox.: Cadastrar/editar resumo, substância, especificação e lote}
\textbf{Contrato:} UI-05. Estoque > Novo Reagente > reutilizar/criar resumo > composição/especificação > lote opcional. Validar identidade e unidades, persistir cada etapa com IDs retornados; densidade confirmada exige nova especificação. Não duplicar resumo por concentração.

\subsubsection{ALM-02 --- Gestor Almox.: Cadastrar frasco fechado ou aberto}
\textbf{Contrato:} UI-06. Estoque > Adicionar Frasco > especificação > almoxarifado/local/lote (opcional) > condição inicial e pesos > validade/destino > revisar > salvar. Transação gera LCQUI-N e grava frasco/histórico; código retornado guia etiqueta física. Repetição da requisição não gera segundo frasco.

\subsubsection{ALM-03 --- Gestor Almox.: Pesquisar, editar localização, abrir e ajustar estoque por pesagem}
\textbf{Contrato:} UI-04 e UI-06. Estoque > filtros > frasco > ação. Abrir recalcula validade; editar localização mantém identidade; pesar exige leitura e motivo, produz histórico de ajuste/evaporação. Conflito com empréstimo obriga recarregar; não sobrescrever peso inicial.

\subsubsection{ALM-04 --- Gestor Almox.: Registrar retirada e autorizar uso didático de vencido}
\textbf{Contrato:} UI-07. Registrar Retirada > código > Professor/Bolsista > local/peso/prazo/finalidade > revisão. Vencido só segue com autorização registrada e ciência explícita; pesquisa excepcional exige justificativa e TCR rastreável Q04. Transação valida disponibilidade e escopo, cria empréstimo e histórico; concorrente perdedor não retira.

\subsubsection{ALM-05 --- Gestor Almox.: Receber devolução de empréstimo próprio ou de outro gestor}
\textbf{Contrato:} UI-07. Registrar Devolução > frasco/pendência > peso/quem devolveu > destino se vencido ou de validade desconhecida > confirmar. O servidor lê o vencimento persistido do frasco (sem recalcular pelo relógio), classifica o retorno (venceu durante, já estava vencido ou validade desconhecida) e calcula consumo e atraso, encerra empréstimo e atualiza frasco/histórico. Peso incompatível é preservado, conclui devolução física com anomalia e consumo pendente para revisão; qualquer gestor vinculado pode receber.

\subsubsection{ALM-06 --- Gestor Almox.: Tratar alertas, quarentena, extravio, reencontro, vazio, quebra e descarte}
\textbf{Contrato:} UI-07, UI-14 e UI-15. Alertas > categoria > frasco > ação/motivo > confirmar. Validar transição e empréstimo ativo; gravar evento, manter identidade e invalidar saldo/alerta. Descarte é terminal; política de remoção lógica em Q07. Em caso de extravio, o frasco preserva seu estado físico, assume \texttt{situacao\_localizacao = EXTRAVIADO} e fica indisponível e preserva o valor anterior de \texttt{saldo\_desconhecido}, o vencimento e \texttt{abertura\_historica\_desconhecida}; um eventual empréstimo ativo sofre encerramento extraordinário. A autorização corrente de descarte é revogada no extravio e a decisão anterior fica no histórico. Em caso de reencontro, o gestor informa o estado físico constatado (\texttt{ABERTO}, \texttt{FECHADO}, \texttt{VAZIO} ou \texttt{QUEBRADO}), validado contra o último estado conhecido antes do extravio: \texttt{FECHADO} só se antes era \texttt{FECHADO} e sem abertura conhecida; \texttt{ABERTO} se antes era \texttt{ABERTO}/\texttt{FECHADO} (marcando abertura histórica desconhecida no segundo caso); \texttt{VAZIO} só se antes já era \texttt{VAZIO}; \texttt{QUEBRADO} representa quebra constatada. Transições impossíveis são rejeitadas. O frasco fica obrigatoriamente com \textit{em\_quarentena = true} e \texttt{INDISPONIVEL}; o saldo permanece como estava, salvo medição ou evento válido, e o empréstimo encerrado não é reaberto. Depois do reencontro, o gestor decide a saída da quarentena (``VOLTAR A DISPONIVEL'' ou ``PENDENTE DE DESCARTE'', com motivo/laudo) e a saída não renova nem recalcula validade; ``VOLTAR A DISPONIVEL'' nunca é oferecido para frasco \texttt{VAZIO} ou \texttt{QUEBRADO}. ``PENDENTE DE DESCARTE'' gera autorização operacional estruturada e mantém o frasco INDISPONIVEL; o descarte só ocorre após essa decisão, nunca diretamente da quarentena. Erros administrativos suportados, como retirada atribuída ao usuário errado, usam ``Corrigir operação'', sem apagar o lançamento original.

\subsubsection{ALM-07 --- Gestor Almox.: Imprimir etiquetas virgens ou segunda via}
\textbf{Contrato:} UI-08. Mais > Imprimir > aba > intervalo/célula ou 1--10 frascos > Gerar > Baixar. Conferir código, imprimir em 100% e colar somente após confirmar cadastro. Virgens não reservam IDs; segunda via audita cada frasco e produz uma etiqueta por ficha.

\subsubsection{ALM-08 --- Gestor Almox.: Gerar relatório de um ou todos os almoxarifados autorizados}
\textbf{Contrato:} UI-12. Relatórios > mensal/período > escopo/filtros > revisar > Gerar/Baixar. Servidor valida 31 dias e permissões, separa g/ml e usa fontes históricas adequadas; erro não anuncia PDF pronto.

\subsubsection{PAT-01 --- Gestor Patr.: Criar resumo, cadastrar e editar bem}
\textbf{Contrato:} UI-09. Patrimônio > Mais > Novo Resumo/Bem, ou detalhe > Editar > preencher/revisar > Salvar. Normalizar a plaqueta por \texttt{trim().toUpperCase()}, validar chave permanente, local, foto e responsável; alteração de nome global exige ação no resumo, reclassificação individual altera referência. Transação cria histórico de cadastro/edição e preserva antes/depois; fan-out derivado não fabrica evento nem conflito de versão.

\subsubsection{PAT-02 --- Gestor Patr.: Pesquisar, filtrar e consultar histórico}
\textbf{Contrato:} UI-04. Patrimônio > prédio/status/letra e filtros opcionais > resumo > plaqueta > detalhe/histórico. Consultar itens e eventos por cursor; mostrar conservação e status distintos, sem substituir histórico pela localização atual.

\subsubsection{PAT-03 --- Gestor Patr.: Analisar, aprovar ou rejeitar pedido}
\textbf{Contrato:} UI-09. Fila/Sino > requisição pendente > comparação > completar dados obrigatórios > justificativa > Aprovar/Rejeitar. Transação revalida M9, confere lock proprietário, versão, chave única, referências e arquivo, persiste decisão, altera/cria bem e histórico se aprovado e libera somente o lock correspondente; notifica solicitante de modo idempotente. Segundo respondente recebe resultado/receipt M7, sem dupla aprovação.

\subsubsection{PAT-04 --- Gestor Patr.: Gerar relação de inservíveis e registrar baixa após SEI}
\textbf{Contrato:} UI-09 e UI-12. Alertas/Relatórios > Inservíveis > filtros > PDF para procedimento externo. Após conclusão no SEI: bem Inservivel > Baixa > upload comprovante > confirmar. Somente confirmação server-side com PDF binariamente validado altera status; histórico e arquivo permanecem disponíveis, enquanto o Bem e sua plaqueta não são excluídos nem reutilizados.

\subsubsection{PAT-05 --- Gestor Patr.: Gerar relatório mensal, personalizado ou individual}
\textbf{Contrato:} UI-12. Relatórios ou detalhe do bem > tipo/período/filtros > Gerar > Baixar/Imprimir. Validar intervalo e datas; documento identifica filtros e instante. Mês em aberto é parcial; erro ou dados ausentes têm retorno distinto.

\paragraph{Q04 e Q14: contrato de aceite e autoatendimento.}
Para PESQUISA\_TCC\_POS ou ESTUDO\_DEGRADACAO\_RESIDUOS com frasco vencido ou validade desconhecida liberado com termo, obter aceite em sessão autenticada do retirante. Persistir tcr\_versao, tcr\_aceito\_em (servidor), tcr\_aceito\_por e tcr\_auditoria\_id no empréstimo, com evento correspondente em Registro\_de\_Auditoria; exigir justificativa\_metodologica de 20--2000 caracteres. Modal/checkbox são UX: backend valida versão vigente, identidade e vínculo do aceite ao frasco, finalidade e operação antes de confirmar a retirada. O TCR registra ciência e responsabilidade e não substitui regras institucionais de segurança química. Autoatendimento só quando não houver outro gestor ativo, calculado pelo servidor com justificativa e notificação à chefia. Registrar abertura\_historica\_desconhecida quando a data anterior não for conhecida; nunca inventar a data.


\paragraph{Exceção metrológica na devolução (HQ-M2-005).}
Se a leitura real contradiz a tara armazenada, o gestor informa exatamente a leitura da balança. O retorno físico encerra a custódia como DEVOLVIDO\_COM\_ANOMALIA, preserva a medição e bloqueia nova retirada; consumo fica pendente até resolução auditada. Referência teórica e tara real são distintas. Confirmar recipiente vazio não autoriza substituir silenciosamente tara real anterior divergente. As opções são nova pesagem, inspeção com confirmação de esgotamento ou recalibração de recipiente vazio conforme o contrato da Seção 10.

\paragraph{Emenda pré-M8: conhecimento metrológico na quebra e no descarte.}
Quebra e descarte preservam \texttt{saldo\_desconhecido} e o conhecimento quantitativo anterior, inclusive no reencontro como \texttt{QUEBRADO}. Não fabricam saldo zero, tara ou pesagem. A baixa operacional não mede a quantidade anteriormente existente. Somente vazio efetivamente confirmado pela rota metrológica válida autoriza saldo zero conhecido. Quarentena, autorização de descarte e terminalidade de \texttt{DESCARTADO} permanecem inalteradas.
